\section{Pendahuluan}
\label{sec:pendahuluan}

Spektroskopi Raman merupakan teknik spektroskopi hamburan inelastik foton yang memberikan informasi mengenai struktur vibrasi, rotasi, serta mode elektronik berfrekuensi rendah dalam suatu sistem materi \cite{raman1928new,long2002raman}. Berbeda dengan spektroskopi serapan inframerah (\textit{Infrared Absorption Spectroscopy}) yang bergantung pada perubahan momen dipol listrik permanen molekul $(\partial \boldsymbol{\mu}/\partial Q \neq 0)$, hamburan Raman didasari oleh modulasi polarizabilitas elektron oleh vibrasi kisi atau koordinat normal molekul $(\partial \boldsymbol{\alpha}/\partial Q \neq 0)$ \cite{bernath2020spectra}. Sifat komplementer ini menjadikan spektroskopi Raman sebagai instrumen diagnostik tak-merusak (\textit{non-destructive}) yang sangat penting dalam sains material kontemporer, khususnya pada eksplorasi struktur berskala nano.

Dalam ranah nanomaterial karbon berdimensi rendah, spektroskopi Raman telah berkembang menjadi standar emas karakterisasi struktural \cite{ferrari2006raman,dresselhaus2005raman,malard2009raman}. Tiga sistem yang menjadi representasi penting dalam fisika zat terkondensasi modern adalah:
\begin{enumerate}
    \item \textbf{Grafena (\textit{Graphene})}: material dua dimensi ($2\text{D}$) berstruktur sarang lebah (\textit{honeycomb}) atom karbon hibridisasi $sp^2$. Spektroskopi Raman mampu membedakan secara tegas jumlah lapisan atomik (monolayer, bilayer, hingga multilayer) melalui dekonvolusi profil garis pita $2D$ ($\sim 2685\text{ cm}^{-1}$), serta mendeteksi kerapatan defek kisi struktural melalui rasio intensitas pita $D$ terhadap pita $G$ ($I_D/I_G$) \cite{ferrari2006raman,cancado2011quantifying}.
    \item \textbf{Tabung Karbon Nano Berdinding Tunggal (\textit{Single-Walled Carbon Nanotubes}, SWCNTs)}: material semu-satu-dimensi ($1\text{D}$) yang dibentuk dari gulungan lembaran grafena sepanjang vektor kiral $\mathbf{C}_h = n\mathbf{a}_1 + m\mathbf{a}_2$. Karakteristik dimensi $1\text{D}$ ini memunculkan mode pernapasan radial (\textit{Radial Breathing Mode}, RBM) pada frekuensi rendah ($100 - 350\text{ cm}^{-1}$) yang berelasi terbalik dengan diameter tabung ($d_t$), serta pemisahan pita $G$ menjadi mode $G^+$ (aksial) dan $G^-$ (sirkumferensial) yang mengidentifikasi sifat elektronik logam versus semikonduktor melalui profil asimetris Breit--Wigner--Fano (BWF) \cite{dresselhaus2005raman}.
    \item \textbf{Heterostruktur van der Waals Grafena/h-BN}: perpaduan lapisan atomik grafena di atas kristal heksagonal boron nitrida (h-BN) yang merupakan isolator celah-pita lebar ($E_g \approx 5.9\text{ eV}$). Substrat h-BN yang bebas dari ikatan menggantung (\textit{dangling bonds}) dan muatan terperangkap menekan fluktuasi potensial lokal (\textit{charge puddles}), menghasilkan penyempitan lebar penuh pada setengah maksimum (\textit{Full Width at Half Maximum}, FWHM) pita $2D$ hingga mendekati batas intrinsik kristal grafena terisolasi \cite{dean2010boron,gorbachev2011hunting}.
\end{enumerate}

Meskipun telaah spektroskopi Raman pada nanomaterial karbon telah banyak dipublikasikan, penjabaran terpadu yang menjembatani penurunan analitik landasan kuantum hamburan inelastik, kaidah seleksi berbasis teori grup, formalisme resonansi Kramers--Heisenberg--Dirac (KHD), hingga pemodelan komputasi berbasis fisika matematis dan verifikasi empiris terkalibrasi masih jarang disajikan dalam satu kerangka kerja yang utuh.

Makalah ini bertujuan untuk membangun jembatan komprehensif tersebut. Pendekatan yang digunakan memadukan dua jalur yang saling memvalidasi:
\begin{enumerate}
    \item \textbf{Jalur Teoretis Analitik}: menurunkan relasi klasik momen dipol terinduksi, formulasi elemen matriks operator koordinat kuantum, tabel karakter teori grup molekular ($D_{\infty h}$, $C_{2v}$, dan $D_{6h}$), teori resonansi KHD, serta persamaan difraksi kisi spektrometer Czerny--Turner.
    \item \textbf{Jalur Komputasi Reproducible}: mengaudit dataset spektrum Raman eksperimental riil (grafena monolayer, bilayer, grafena terdefek, SWCNT logam, SWCNT semikonduktor, dan heterostruktur grafena/h-BN), menerapkan pipeline pemrosesan sinyal terkontrol (eliminasi \textit{cosmic ray}, koreksi garis dasar \textit{Asymmetric Least Squares}, dan penghalusan Savitzky--Golay), melakukan dekonvolusi profil garis spektral non-linear, menyusun simulasi fisika maju (\textit{forward physical modeling}) dari parameter molekular, serta mengevaluasi deviasi residual secara kuantitatif melalui metrik MAE, RMSE, Pearson $r$, dan koefisien determinasi $R^2$.
\end{enumerate}

Sistematika makalah ini disusun sebagai berikut: Bagian \ref{sec:dasar_raman} menjabarkan dasar teori hamburan Raman secara klasik dan kuantum. Bagian \ref{sec:kaidah_seleksi} menurunkan kaidah seleksi vibrasi dan rotasi serta analisis teori grup. Bagian \ref{sec:resonance_raman} membahas teori \textit{resonance Raman} dan formalisme KHD. Bagian \ref{sec:instrumentasi} menguraikan arsitektur optik dan detektor spektrometer mikro-Raman. Bagian \ref{sec:analisis_simulasi} menyajikan analisis data eksperimen, simulasi maju, dan perbandingan kuantitatif keduanya. Bagian \ref{sec:kesimpulan} menutup naskah dengan kesimpulan serta tinjauan keterbatasan model.
